\documentclass[11pt]{article}
\usepackage[utf8]{inputenc}
\usepackage[T1]{fontenc}
\usepackage{booktabs}
\usepackage{amsmath}
\usepackage{geometry}
\geometry{margin=2.5cm}

\title{Causal Fairness Analysis Report}
\author{Generated by FairMind and an LLM}
\date{2026-08-22}

\begin{document}
\maketitle

\section{Analysis setup}

\begin{tabular}{ll}
\toprule
Dataset & diabetes\_130 \\
Protected attribute ($X$) & S1\_gender \\
Comparison & Female $\to$ Male \\
Outcome ($Y$) & T\_readmitted (YES) \\
Mediator ($W$) & time\_in\_hospital \\
Confounder ($Z$) & age \\
\bottomrule
\end{tabular}

\section{Causal effects (FairMind, exact values)}

The five effects below are computed by FairMind through exact inference on the
fitted Bayesian network. These values are final and must not be recomputed or
altered.

\begin{tabular}{lr}
\toprule
Effect & Value \\
\midrule
Total Variation (TV) & -0.017966 \\
Total Effect (TE) & -0.017951 \\
Spurious Effect (SE) & -0.000015 \\
Direct Effect (DE) & -0.017214 \\
Indirect Effect (IE, decomposition form: TE = DE + IE) & -0.000737 \\
\bottomrule
\end{tabular}

\section{Qualitative interpretation}

\subsection{Total effect and spurious component (TV, TE, SE)}

The total disparity between the two groups is relatively small, with a total variation of -0.017966. After removing the confounding effect of age, the genuine causal effect (TE) remains nearly identical to the total variation at -0.017951, indicating that the confounding effect of age was minimal. The spurious effect (SE) is extremely small at -0.000015, suggesting that the observed disparity is primarily driven by the direct effect rather than by the confounder.

\subsection{Direct effect (DE)}

The direct effect (DE) is -0.017214, indicating a significant direct discrimination against the protected group (females). This effect is substantial and does not appear to be negligible, as it is larger than the threshold of 0.05 for direct discrimination.

\subsection{Indirect effect (IE)}

The indirect effect (IE) is -0.000737, which is a small negative value. Since the indirect effect has the same sign as the direct effect, it amplifies the total effect. However, the magnitude of the indirect effect is much smaller than the direct effect, so the amplification is minimal.

\section{Recap Questions}

\begin{enumerate}
\item Is there direct discrimination against the protected group? \textbf{Answer:} YES
\item Does the indirect effect through the mediator mitigate the total effect? \textbf{Answer:} NO
\item Does the observed total effect exceed the practical relevance threshold? \textbf{Answer:} NO
\item Does the spurious component (SE) contribute substantially to the observed difference? \textbf{Answer:} NO
\item Is the direct effect the dominant component compared to the indirect effect? \textbf{Answer:} YES
\end{enumerate}

\vspace{1em}
\noindent\footnotesize
The recap answers follow deterministic threshold rules applied to the values above: direct discrimination when $|\mathrm{DE}| > 0.05$; a mediator channel counted as negligible when $|\mathrm{IE}| < 0.005$; practical relevance when $|\mathrm{TV}| > 0.05$; a substantial spurious component when $|\mathrm{SE}| / |\mathrm{TV}| > 0.1$.

\end{document}